\subsection{Training Infrastructure}




Efficient training of large-scale autoregressive denoising models like \magi requires carefully tailored distributed training infrastructure. Existing distributed training frameworks, such as Megatron~\citep{shoeybi2020megatronlm} and DeepSpeed~\citep{rajbhandari2020zeromemoryoptimizationstraining} are primarily designed for large language models (LLMs). However, \magi differs significantly from LLMs in both algorithmic side and data side.

On the algorithmic side, \magi integrates both autoregressive and denoising modeling paradigms, resulting in a model architecture that is notably more complex than that of typical LLMs. It incorporates components such as gating, cross-attention, and block-causal attention that are rarely used in language models.


On the data side, a single video training example typically contains tens to hundreds of times more tokens than a text example. Furthermore, ensuring the temporal and semantic integrity of video content imposes strict constraints, making it infeasible to directly apply common data processing strategies from LLMs, such as arbitrary sequence truncation or concatenation of multiple samples into a single training sequence offline, in the context of video generation.

These fundamental differences introduce unique challenges, necessitating a new, purpose-built distributed system design. In this section, we propose novel solutions to address these challenges to enable efficient and scalable training of \magi. 


Specifically, the training of \magi leverages a combination of data parallelism (DP), context parallelism (CP), and tensor parallelism (TP). To address the DP load imbalance caused by variable-length video sequence and the insufficient GPU utilization on short token sequences, we introduce a distributed Packing and Padding (PnP) during training, that performs online batching of video data in each training iteration. This strategy mitigates GPU bubbles thereby significantly improving overall training efficiency (Sec.~\ref{sec:train_infra_pnp}).

Due to the use of PnP and the inherent demands of block-causal attention in \magi, we require an efficient attention implementation capable of supporting highly flexible attention masks. Additionally, given the extremely long token sequences typical in video training data, native support for context parallelism is essential. To address these requirements, we propose MagiAttention: a scalable distributed attention mechanism that efficiently handles diverse attention masks and is optimized for ultra-long sequences (Sec.~\ref{sec:magiattn}).

Through the above innovations, we enable the efficient training of \magi. However, while developing the \magi training system, we identified several limitations in existing large-scale training frameworks. For instance, most current frameworks (including ours) do not treat verifiable numerical accuracy in distributed environments as a first-class design concern. Moreover, the tight coupling between algorithm development and infrastructure implementation often creates friction between algorithm researchers and infrastructure engineers, hindering efficient collaboration. To address these challenges, we discuss potential directions and design principles for next-generation training infrastructure in Sec.~\ref{sec:train_infra_odin}, with the goal of providing insights and practical guidance for the broader research and engineering community.








\subsubsection{Distributed Packing and Padding}\label{sec:train_infra_pnp}



Due to the integrity constraints of video data and the variability in video lengths and resolutions, we adopt a Packing and Padding (PnP) strategy~\citep{sirluk2024efficient, kundu2024enhancingtrainingefficiencyusing} to batch video samples in a way that minimizes excessive padding and reduces unnecessary computational overhead in distributed training scenarios. Moreover, the data composition is frequently adjusted during the training of \magi (See Sec.~\ref{sec:data_adjustment}), and to accommodate such flexibility, we employ an online PnP strategy instead of a offline approach.

The core idea of PnP is to efficiently utilize GPU resources by concatenating multiple short sequences into a batch while minimizing redundant filling. The offline formulation of this problem aligns with the classic bin-packing problem: given a set of input samples, the goal is to pack them into a set of bins, each with a fixed capacity \texttt{max\_length}, while minimizing overall unused space. Although this problem is NP-complete, it can be efficiently approximated in practice using the First-Fit Decreasing (FFD)~\citep{dosa2007tight} greedy algorithm.

In our online setting, we must process streaming data inputs while ensuring compatibility with the 3D parallelism strategy employed during training. To this end, we reformulate the problem as follows: given M candidate samples, we aim to pack them into N bins of size \texttt{max\_length}, minimizing overall space waste. Here, M denotes the size of the candidate pool with M $\gg$ N; N must be divisible by the \texttt{DP\_SIZE}; and \texttt{max\_length} must be divisible by \texttt{TP\_SIZE}$\times$\texttt{CP\_SIZE}.

In practice, we extend the FFD algorithm with custom heuristics to support efficient online packing under these constraints. This approach enables us to achieve a $99\%$ capacity utilization rate and the differences between different DP groups can be neglected, thus substantially reducing computational overhead during training.




































\subsubsection{MagiAttention: Towards Linear Scalability for Ultra-Long and Heterogeneous Mask Training.}\label{sec:magiattn}

\begin{figure}[!htbp]
    \centering
    \includegraphics[width=\linewidth]{infra/train_infra/figs/magiattn_overview_v3.pdf}
    \caption{Overview of MagiAttention: (1) Flex-Flash-Attention(FFA), an efficient attention supports flexible mask patterns and native considers distribution requirements; (2) The \textit{dispatch solver} shards and dispatches packed data with ultra-long contexts and heterogeneous masks, ensuring load-balanced computation; (3) Group-Cast and Group-Reduce primitives eliminate redundant communication; (4) The \textit{adaptive multi-stage overlap} strategy effectively hides communication latency; (5) Forward and backward timelines of MagiAttention. With all techniques together, MagiAttention reach linear scalability under diverse scenarios.}

    
    \label{fig:magiattn_overview}
\end{figure}


Training large-scale autoregressive diffusion models like \magi for video generation presents two major challenges: 

\begin{itemize}
    \item The extremely long context length of video tokens, which reaching up to 4 million during training, results in prohibitive computational and memory overhead. Context-Parallelism (CP) is designed for dealing such long context challenge, but existing state-of-the-art CP methods~\citep{jacobs2023deepspeed, liu2023ringattentionblockwisetransformers, fang2024uspunifiedsequenceparallelism, gu2024loongtrainefficienttraininglongsequence, chen2024longvilascalinglongcontextvisual} face scalability limitations that face scalability limitations due to size constraints or the high communication overhead inherent in inefficient ring-style point-to-point (P2P) patterns. While recent efforts~\citep{wang2024datacentricheterogeneityadaptivesequenceparallelism, zhang2024dcp, ge2025bytescaleefficientscalingllm} dynamically adjust CP sizes to avoid unnecessary sharding and redundant communication for shorter sequences, they still incur extra memory overhead for multiple NCCL process groups and involve complex scheduling to balance loads and synchronize across different subsets of ranks.
    \item The combination of block-causal attention and Packing-and-Padding introduces highly complex attention mask patterns (Sec.\ref{sec:train_infra_pnp}), which cannot be efficiently handled by existing attention implementations.
\end{itemize}





To address the aforementioned challenges, we propose MagiAttention, which aims to support a wide variety of attention mask types (\emph{i.e.}, kernel flexibility) while achieving linear scalability with respect to context-parallel (CP) size across a broad range of scenarios. Achieving this goal depends on meeting the following fundamental conditions:

\begin{itemize}
    \item \emph{Linearly Scalable Attention Kernel}: The performance of the attention kernel should not degradate as CP size increases. To this end, we introduce \emph{Flex-Flash-Attention}, an extension of FlashAttention-3 (FA3), which native considers the efficiency impact of attention mask partitioning in distributed environments. It supports distributable mask representations with a tailored kernel implementation to ensure scalability while accommodating a broader range of attention mask types.
    \item \emph{Balanced Computational Workloads}: Imbalances in the computational load across CP ranks lead to unavoidable idle bubbles that hinder scalability. MagiAttention is natively designed to ensure \emph{Computation Load Balancing}, mitigating such inefficiencies.
    \item \emph{Full Overlap of Communication and Computation}: Without sufficient overlap, increasing CP size results in communication-induced idle time on GPUs, impairing scalability. MagiAttention introduces novel \emph{Zero-Redundant Communication Primitives} to minimize communication overhead, along with an \emph{Adaptive Multi-Stage Overlap} strategy that enables effective communication-computation overlap.
\end{itemize}

The overview of MagiAttention is shown in Fig.~\ref{fig:magiattn_overview}, and we will introduce key designs in the following, with comprehensive experimental results presented in Appendix~\ref{appendix:magiattn_exps}.



















\paragraph{Flex-Flash-Attention.}\label{sec:magiattn_ffa}
\input{infra/train_infra/figs/mask_with_attn_slice}

FlashAttention~\citep{dao2022flashattention, dao2023flashattention, shah2024flashattention3fastaccurateattention} is a foundational technique in large-scale model training for its superior performance and support for varlen-packed data with causal attention masks. However, it offers limited support for irregular attention masks, particularly when such patterns are distributed across CP ranks, resulting in increased complexity and underscoring the need for a more flexible attention kernel~\citep{pytorch_sdpa, dong2024flexattentionprogrammingmodel,wang2025flashmaskefficientrichmask} without compromising performance.

Therefore, we introduce Flex-Flash-Attention (FFA), which is natively designed for distribution scenarios and provides greater flexibility in handling diverse attention mask types. The core idea behind FFA is to generalize a distributable formulation for irregular attention masks by decomposing the entire mask into multiple computational units, each referred to as an \texttt{AttnSlice}. Each \texttt{AttnSlice} is defined by a triplet \texttt{\(QRange, KRange, MaskType\)}, which specifies a submask with a basic shape bounded by a contiguous 2D query-key region (see Fig.~\ref{fig:attnslice_interpret}). Using this formulation, a wide variety of commonly used attention masks (Fig.~\ref{fig:mask_with_attn_slice}) (including our varlen block-causal mask) can be expressed as a composition of multiple such triplets, making FFA highly suitable for distributed attention computation.


Built on FA3 kernels, Flex-Flash-Attention leverages NVIDIA Hopper GPUs' TMA feature~\citep{nvidia2024accelerating} and introduces slice-level parallelism with atomic operations for correctness (Fig~\ref{fig:ffa_slice_parallel}), achieving comparable MFU to FA3 while supporting the flexible \texttt{AttnSlice} formulation~\footnote{Redundant computation from padding tokens is excluded by easily passing empty \texttt{QRange} or \texttt{KRange}.} (see Appendix~\ref{appendix:magiattn_exps} for benchmarks).








\paragraph{Computation Load-Balance.}\label{sec:magiattn_comp}




In context-parallelism (CP) settings, different CP ranks may be assigned heterogeneous attention masks, resulting in imbalanced computational workloads across ranks. Ring-Attention~\citep{ring_flash_attention_issue2} employs a specialized partitioning strategy designed specifically for causal attention, which limits its applicability to more general attention patterns. To overcome this limitation, we propose a generic and efficient \texttt{dispatch solver} that enables balanced workload distribution across CP ranks for a broad range of attention types.

First, to enable finer-grained control, we propose a chunk-wise permutable sharding strategy (Fig~\ref{fig:magiattn_overview} (2)). Specifically, the entire mask is evenly partitioned along the query-dimension into dispatch chunks, each associated with a submask area: $\{(C_i, \mathrm{Area}(C_i))\}_{i=1}^n$, where $C_i$ indicates i-th dispatch chunk, $\mathrm{Area}(C_i)$ is the mask area of $C_i$, $n$ is $\frac{seqlen}{dispatch\_chunk\_size}$, and $dispatch\_chunk\_size$ is a hyperparameter controlling granularity. 
These dispatch chunks are then equally assigned to $cp\_size$ buckets, with each bucket containing the exact same number of dispatch chunks to ensure token-level load balance in non-attention modules, attaching with a summed submask area, denoted as $\{(B_j, \mathrm{SumArea}(B_j))\}_{j=1}^{cp\_size}$.


With above strategy, we could fine-grained control the computational workloads of each CP rank, and the load-balancing dispatch becomes a combinatorial optimization problem, defined as finding an optimal mapping function $f^*: \{C_i\}_{i=1}^n \rightarrow \{B_j\}_{j=1}^{cp\_size}$ as follows

\begin{align}
    &f^* = \arg \min\limits_{f}\max\limits_{j}\left\{\mathrm{SumArea}(B_j)\right\} \label{eq:comp_load_balance}\\
    &\text{s.t.}\;\;|B_j| = \frac{n}{cp\_size}, \;\; seqlen \;\%\; (cp\_size \times dispatch\_chunk\_size) = 0\nonumber
\end{align}


However, this optimization is a known NP-hard problem, making it impractical to find an optimal solution on-the-fly during each training iteration, especially given the varying mask patterns across micro-batches. Thus, we propose an efficient greedy algorithm (as shown in Alg.~\ref{alg:minhp}) that provides a suboptimal yet effective solution within $O(n\log n)$ complexity.













\paragraph{Zero-Redundant Communication Primitives.}\label{sec:magiattn_comm}



The existing ring-style implementation uses point-to-point \texttt{send}/\texttt{recv} communication primitives, which cannot provide sufficient communication granularity, resulting in redundant communication. Take causal mask as an example, we analyze the redundant communication by recording the distribution of remote key-value ($\mathrm{KV}$) requests and their gradients ($\mathrm{dKV}$) under sparse attention masks. As shown in Fig~\ref{fig:ring-p2p-redundancy}, $\mathrm{KV}_0$ is required by all queries and should be sent to all devices via Broad-Cast in the forward pass, with $\mathrm{dKV}_0$ reduced via All-Reduce in the backward pass. In contrast, $\mathrm{KV}_7$ is only needed by its host device but still circulates through all devices, and this redundancy intensifies in varlen scenarios.



To address this, we introduce two communication primitives: \texttt{group-cast} and \texttt{group-reduce}, which model the communication patterns of low-demand $\mathrm{KV}$ and $\mathrm{dKV}$ (Fig~\ref{fig:group-gather-reduce-all2allv}). For example, in the causal mask, $\mathrm{KV}_5$ on $\mathrm{rank}_2$ is required only by $\{\mathrm{Q}_6,\mathrm{Q}_7\}$ and should be sent exclusively to the target ranks $\{\mathrm{rank}_0, \mathrm{rank}_1\}$ via \texttt{group-cast}, while the partial $\mathrm{dKV}_5$ is collected and reduced back to $\mathrm{rank}_2$ via \texttt{group-reduce} accordingly.

As no existing communication kernels support these primitives, we prototype them using \texttt{all-to-all-v} (Fig~\ref{fig:group-gather-reduce-all2allv}), achieving zero-redundant communication in both forward and backward passes. However, this approach introduces extra pre-/post-processing overhead, similar to (un)permutation in expert parallelism (EP)~\citep{gale2022megablocks}. While kernel fusion mitigates the overhead, a dedicated implementation of \texttt{group-cast} and \texttt{group-reduce} remains a key direction for future work.











\paragraph{Adaptive Multi-Stage Overlap.}\label{sec:magiattn_overlap}



Leveraging previous optimizations, we achieve high-performance computation through an efficient kernel and balanced workload dispatch, while minimizing communication overhead with our new primitives. To drive true linear scalability, we further improve end-to-end performance by introducing a multi-stage compute-communication overlap strategy, that effectively hides communication latency and adaptively optimizes overlap through manual or automatic tuning.

Similar to prior works~\citep{liu2023ringattentionblockwisetransformers,zhao2023pytorch,async_tensor_parallelism_in_pytorch}, we schedule pipeline stages to overlap computation with communication for both forward and backward passes (Fig~\ref{fig:multi_stage_overlap_fwd_bwd}). Each $\mathrm{rank}_i$ first partitions its remote $\mathrm{KV}$/$\mathrm{dKV}$ communication into stages. 
In the forward pass, the scheduler first launches the \texttt{group-cast} kernel to prefetch the next remote $\mathrm{KV}$, then asynchronously executes the FFA kernel for partial attention computation, hiding all communication behind computation~\footnote{To prevent all SMs from being occupied by the attention kernel, we ensure the communication kernel picked first by setting \texttt{CUDA\_DEVICE\_MAX\_CONNECTIONS}=1~\citep{cuda_device_max_connections_issue}.}. 
In the backward pass, besides prefetching the next $\mathrm{KV}$, the \texttt{group-reduce} kernel reduces the last $\mathrm{dKV}$ in a separate CUDA stream before launching the FFA kernel for the current stage, ensuring communication is overlapped across all stages except the final $\mathrm{dKV}$ reduction~\footnote{Due to PyTorch's one-to-one mapping for process groups and collective communication streams including \texttt{all-to-all-v}~\citep{collectives_nccl_stream_issue}, we internally use an additional CP group for \texttt{group-reduce} to enable full overlap between communication kernels in the backward pass.}.

To adaptively control overlap granularity, we further introduce a tunable hyperparameter, $num\_stages$, accounting for varying compute-to-communication ratios across training setups, microbatches, or between forward and backward passes. This parameter can be manually configured or automatically determined by our \textit{overlap solver}, with a simple dynamic search algorithm (See Alg.~\ref{alg:dynamic-overlap-search} for more details).























    
    
    












\subsubsection{Rethinking System Design for Robust Distributed Training Frameworks with DTensor}\label{sec:train_infra_odin}

As large-scale models continue to evolve, the growing  complexity of training procedures has exposed fundamental limitations in existing distributed training frameworks~\citep{shoeybi2020megatronlm, rajbhandari2020zeromemoryoptimizationstraining}. Two major bottlenecks are particularly prominent:
\begin{itemize}
    \item Lack of testability by design. Most frameworks were not initially built with testability as a first-class feature, resulting in fragile infrastructure with limited maintainability and reliability;
    \item Tight coupling between model implementation and parallelization strategy. This entanglement prevents algorithm researchers and system engineers from working independently, hindering collaboration and modular development
\end{itemize}
We argue that next-generation distributed training frameworks must directly address these two pain points to support large-scale model research and deployment.

Inspired by early explorations~\citep{xu2021gspmdgeneralscalableparallelization, yuan2022oneflowredesigndistributeddeep} and PyTorch’s pioneering implementations~\citep{zhao2023pytorch, pytorchdtensor2024, liang2024torchtitanonestoppytorchnative}, we propose a blueprint for redesigning robust distributed training frameworks based on Pytorch Distributed Tensor (DTensor)~\citep{pytorchdtensor2024} and Parallel Plan:

\paragraph{DTensor} PyTorch DTensor introduces three parallel placements: \texttt{Replicated}, \texttt{Shard}, and \texttt{Partial}, alongside a distributed initialization strategy to maintain placement semantics~\citep{pytorch_offset_rngtracker_2025}, and a propagation mechanism that deduces output placements from input ones for supported ops, triggering communication as needed~\footnote{In practice, DTensor selects communication patterns based on estimated redistribution cost, but these estimates are often inaccurate.}~\citep{pytorch_sharding_prop_2025}. While it supports basic ops including naive distributed matmul, its current implementations lack the generality to handle more complex yet commonly scenarios in modern training workflows, as shown in Tab.~\ref{tab:dtensor_placements}.

\paragraph{Parallel Plan} 
Parallel Plan provides a declarative interface for specifying parallelization strategies across model submodules. It works in conjunction with the \texttt{parallelize\_module} function and is built on top of DTensor. However, its current capabilities are mostly limited to tensor parallelism (TP) and do not generalize well to other parallelism.

In our architecture design, we extend both DTensor and Parallel Plan to support a broader range of usages. These extensions enable the following key features:

\paragraph{Decoupling Modeling from Parallelization.} This feature allows model researchers to concentrate on model design and algorithm development without needing to manage low-level parallelism details. At the same time, infrastructure engineers can independently optimize parallelization strategies without modifying model implementation. This clear separation of concerns enables more efficient collaboration and improved training throughput.


\paragraph{High-Precision Alignment with Non-Distributed Oracles.} 
By disabling all parallel plans, we can seamlessly revert to non-distributed configurations, yielding "pure" model code that serves as a baseline or oracle for evaluating distributed correctness. To ensure alignment within a relative error of $10^{-8}$, we upcast tensors to higher precision~\footnote{In our experiments, \texttt{float32} is insufficient for fully alignment; \texttt{float64} suffices in most cases.}, enforce deterministic algorithms~\citep{pytorch2024reproducibility}, and control randomness using consistent seed management. This design enables precise infrastructure testing, ultimately improving reliability and debuggability.










